\section{Empirical Results}
\label{sec:results}

This section evaluates the two empirical channels defined previously. The first asks whether Variance Risk Premium variables improve dynamic equity--bond allocation. The second evaluates the model-based direct variance-payoff approximation. These results are interpreted jointly, but they are not mechanically ranked against one another because the two layers use different payoff mappings and different aligned samples.

The allocation comparison contains 184 out-of-sample observations in the United States and 122 in Europe. The direct-variance strategy comparison contains 232 observations in the United States and 170 in Europe. The underlying payoff diagnostics contain 256 settlement observations in the United States and 194 in Europe. All comparisons below therefore preserve the sample appropriate to the evidence layer being discussed.

For descriptive comparison, the reported model-family leaders are selected ex post from the evaluated out-of-sample specification set. Their ranking should therefore not be interpreted as the outcome of a pre-specified model-selection rule.

Figure~\ref{fig:vrp-components} provides descriptive context for the variance measures used throughout the empirical analysis. The panels display the implied- and realized-variance components underlying the empirical VRP construction in the two markets. The purpose of the figure is descriptive rather than causal: it illustrates the time variation of the variance measures that subsequently enter the allocation and direct-payoff exercises.

\begin{figure}[htbp]
\centering

\begin{subfigure}[t]{0.49\textwidth}
\centering
\includegraphics[
    width=\linewidth
]{us_mvp1_vrp_components.png}
\caption{United States}
\end{subfigure}
\hfill
\begin{subfigure}[t]{0.49\textwidth}
\centering
\includegraphics[
    width=\linewidth
]{eu_mvp1_vrp_components.png}
\caption{Europe}
\end{subfigure}

\caption{Implied variance, realized variance, and empirical VRP components}
\label{fig:vrp-components}
\end{figure}

Before turning to the summary statistics, Figure~\ref{fig:allocation-cumulative} displays the cumulative paths of the final implementable allocation comparisons in both markets. Presenting the strategy paths first provides a visual description of relative performance and drawdown behavior before the numerical tables are considered.

\begin{figure}[htbp]
\centering

\begin{subfigure}[t]{0.49\textwidth}
\centering
\includegraphics[
    width=\linewidth
]{us_final_implementable_cumulative_returns.png}
\caption{United States}
\end{subfigure}
\hfill
\begin{subfigure}[t]{0.49\textwidth}
\centering
\includegraphics[
    width=\linewidth
]{eu_final_implementable_cumulative_returns.png}
\caption{Europe}
\end{subfigure}

\caption{Cumulative performance of the final implementable allocation comparisons}
\label{fig:allocation-cumulative}
\end{figure}


Figure~\ref{fig:historical-equity-indices} provides complementary market context by showing
the cumulative evolution of the equity components used in the allocation exercise. Each series
is normalized to 100 at the beginning of its corresponding out-of-sample allocation period.
The figure is descriptive: because the United States and European allocation samples begin at
different dates, the two normalized paths should not be interpreted as a common-period
performance comparison.

\begin{figure}[htbp]
\centering
\includegraphics[
    width=0.88\textwidth
]{historical_equity_indices.png}
\caption{Historical evolution of the equity-market components used in the allocation exercise}
\label{fig:historical-equity-indices}
\end{figure}

\subsection{Equity--bond allocation: United States}

Table~\ref{tab:us_allocation} reports the principal United States allocation results. Buy-and-hold equity generates the highest annualized return, at 13.96\%, but also the highest volatility and the largest drawdown among the benchmark portfolios. The 60/40 and equal-weighted portfolios reduce risk substantially. Their Sharpe ratios, 1.024 and 1.025 respectively, provide demanding benchmarks for the dynamic models.

\begin{table}[htbp]
\centering
\small
\caption{United States allocation results}
\label{tab:us_allocation}
\begin{tabularx}{\textwidth}{Xrrrrrr}
\toprule
Strategy
& Return
& Vol.
& Sharpe
& MDD
& CVaR
& Turnover \\
\midrule
Buy-and-Hold Equity
& 13.96\%
& 14.20\%
& 0.995
& -23.93\%
& -8.18\%
& 0.00\% \\
60/40
& 9.39\%
& 9.21\%
& 1.024
& -20.06\%
& -5.20\%
& 1.47\% \\
1/N Equity--Bond
& 8.23\%
& 8.05\%
& 1.025
& -19.10\%
& -4.58\%
& 1.53\% \\
HMM RV + Log VRP
& 8.36\%
& 8.24\%
& 1.019
& -16.67\%
& -4.86\%
& 25.35\% \\
RSM RV + Raw VRP
& 9.61\%
& 10.25\%
& 0.949
& -15.43\%
& -5.53\%
& 36.69\% \\
ML Logistic Base
& 7.58\%
& 7.60\%
& 1.002
& -18.54\%
& -4.35\%
& 17.39\% \\
\bottomrule
\end{tabularx}
\end{table}

The strongest HMM by Sharpe ratio is HMM RV + Log VRP, with a Sharpe ratio of 1.019. This is close to the balanced benchmarks but remains below both 60/40 and equal weighting. Its main improvement is on downside risk: maximum drawdown falls to -16.67\%, compared with -20.06\% for 60/40, -19.10\% for equal weighting, and -23.93\% for buy-and-hold equity.

That improvement is not free. Average turnover rises to 25.35\%, compared with approximately 1.5\% for the static balanced portfolios. The HMM RV + Raw VRP specification produces a higher annualized return of 10.67\% and a Sharpe ratio of 1.014, but does not improve the broad risk-adjusted ranking.

The strongest RSM specification is RSM RV + Raw VRP. Its maximum drawdown of -15.43\% is the smallest among the principal United States allocation strategies, but its Sharpe ratio is only 0.949 and average turnover reaches 36.69\%. The reduction in drawdown is therefore accompanied by substantially greater implementation intensity.

The strongest machine-learning allocation is ML Logistic Base, with a Sharpe ratio of 1.002. Importantly, this leading United States ML specification does not contain the VRP-enhanced feature set. The United States evidence therefore does not support a general claim that adding VRP variables improves allocation performance.

Taken together, the United States allocation results indicate that the dynamic models can alter the distribution of portfolio risk, particularly drawdown, without establishing robust dominance over simple balanced portfolios. The distinction between a better downside statistic and a better investment strategy is economically important.

The implementation mechanism of the selected United States HMM is illustrated in Figure~\ref{fig:hmm-mechanism}. The first panel reports the filtered stress probabilities and the second reports the corresponding equity weights. The figure shows how changes in the inferred latent state are mapped into portfolio exposure rather than treating the HMM as a purely statistical classification device.

\begin{figure}[htbp]
\centering

\begin{subfigure}[t]{0.49\textwidth}
\centering
\includegraphics[
    width=\linewidth
]{us_mvp2_hmm_stress_probabilities.png}
\caption{Filtered stress probabilities}
\end{subfigure}
\hfill
\begin{subfigure}[t]{0.49\textwidth}
\centering
\includegraphics[
    width=\linewidth
]{us_mvp2_hmm_equity_weights.png}
\caption{Resulting equity weights}
\end{subfigure}

\caption{Selected United States HMM: state probability and portfolio mapping}
\label{fig:hmm-mechanism}
\end{figure}

\subsection{Equity--bond allocation: Europe}

The European results are materially weaker in absolute risk-adjusted terms. Table~\ref{tab:eu_allocation} reports the principal comparison.

\begin{table}[htbp]
\centering
\small
\caption{European allocation results}
\label{tab:eu_allocation}
\begin{tabularx}{\textwidth}{Xrrrrrr}
\toprule
Strategy
& Return
& Vol.
& Sharpe
& MDD
& CVaR
& Turnover \\
\midrule
Buy-and-Hold Equity
& 3.98\%
& 17.17\%
& 0.313
& -35.74\%
& -9.86\%
& 0.00\% \\
60/40
& 2.68\%
& 11.24\%
& 0.292
& -20.32\%
& -6.55\%
& 1.74\% \\
1/N Equity--Bond
& 2.30\%
& 9.84\%
& 0.280
& -19.66\%
& -5.85\%
& 1.81\% \\
HMM RV
& 1.80\%
& 10.94\%
& 0.218
& -19.39\%
& -6.20\%
& 17.76\% \\
RSM RV + Raw VRP
& 2.55\%
& 11.24\%
& 0.281
& -20.47\%
& -6.74\%
& 34.00\% \\
ML Random Forest + VRP
& 2.26\%
& 9.99\%
& 0.274
& -19.17\%
& -5.82\%
& 15.38\% \\
\bottomrule
\end{tabularx}
\end{table}

Buy-and-hold equity produces the highest benchmark Sharpe ratio, at 0.313, while the balanced portfolios substantially reduce drawdown. The 60/40 portfolio reaches a Sharpe ratio of 0.292 and equal weighting 0.280.

The strongest HMM is HMM RV, with a Sharpe ratio of 0.218. Adding raw VRP reduces the HMM Sharpe ratio to 0.166, while adding LogVRP produces 0.179. Turnover also increases. Within this model family, the corrected European evidence therefore provides no performance improvement from adding VRP.

The strongest RSM is RSM RV + Raw VRP. Its Sharpe ratio of 0.281 is essentially equal to the equal-weighted benchmark but below both 60/40 and buy-and-hold equity. Average turnover is 34.00\%, compared with less than 2\% for the two balanced benchmarks. Raw VRP therefore contributes useful state information within this specification, but not enough to establish economic dominance.

The strongest European machine-learning allocation is ML Random Forest + VRP. It generates a Sharpe ratio of 0.274 and a maximum drawdown of -19.17\%. The corresponding classifier provides some evidence of incremental predictive information: ROC-AUC is approximately 0.619, precision--recall AUC 0.468, and the Brier score 0.215. However, the portfolio remains below the benchmark Sharpe ratios.

This distinction between prediction and portfolio value is reinforced by the welfare results for the European ML layer. At a risk-aversion coefficient of five, its mean--variance certainty-equivalent advantage is only about 12 basis points relative to 60/40 and approximately -10 basis points relative to equal weighting, with confidence intervals that include zero. Statistical predictability therefore does not translate into robust investor-welfare dominance.

The European allocation conclusion is consequently narrow. VRP-related variables are not empirically absent, and they can contribute information in selected RSM and machine-learning specifications. They do not, however, produce robust superiority in the tested equity--bond allocation framework.

\subsection{Cross-market allocation evidence}

The cross-market allocation evidence rejects a simple interpretation in which greater model complexity automatically creates economic value. In the United States, the strongest dynamic models are close to the balanced benchmarks and sometimes improve drawdown, but they require considerably more turnover. In Europe, the regime and machine-learning models remain below the strongest traditional benchmarks.

The role of VRP is also model dependent. HMM RV + Log VRP is the strongest HMM in the United States, whereas the strongest European HMM excludes VRP. Raw VRP improves the relative ranking of the RSM family in both markets, but neither RSM becomes a dominant allocation strategy. The leading United States ML specification excludes VRP, while the leading European ML specification includes it.

The allocation layer therefore provides evidence that variance information can affect estimated states and selected portfolio-risk characteristics, but it does not establish a universal allocation premium.



\subsection{Underlying direct variance-payoff diagnostics}

The second evidence layer begins with the modeled short-variance payoff itself. Table~\ref{tab:payoff_diag} summarizes the underlying settlement observations before the final strategy alignment.

\begin{table}[htbp]
\centering
\small
\caption{Underlying short-variance payoff diagnostics}
\label{tab:payoff_diag}
\begin{tabularx}{\textwidth}{Xrrrrrr}
\toprule
Market
& Obs.
& Positive
& Mean norm.
& Skew
& 1\% quantile
& Worst \\
\midrule
United States
& 256
& 82.42\%
& 0.253
& -3.849
& -4.080
& -4.432 \\
Europe
& 194
& 78.35\%
& 0.239
& -2.231
& -2.060
& -2.439 \\
\bottomrule
\end{tabularx}
\end{table}

The average variance strike is 4.35\% in the United States and 5.22\% in Europe, compared with average realized variance of 3.57\% and 4.02\%, respectively. The corresponding mean short-variance spreads are positive in both markets.

Positive carry is frequent but not symmetric. The normalized payoff is positive in 82.42\% of United States observations and 78.35\% of European observations. At the same time, skewness is strongly negative: -3.849 in the United States and -2.231 in Europe. Approximately 4.69\% of United States observations and 4.12\% of European observations generate normalized losses below -100\%.

These diagnostics are central to interpretation. A high positive-payoff frequency does not imply an arbitrage or a low-risk return process. The short-variance mechanism earns frequent positive carry while remaining exposed to infrequent, disproportionately large losses. Risk targeting, exposure constraints and welfare analysis are therefore necessary before drawing an economic conclusion.

\subsection{United States direct-variance evidence}

The selected United States direct-variance specification is the 10\% annual-volatility target activated by the High-VRP gate. Table~\ref{tab:direct_main} reports the selected direct strategies for both markets.

\begin{table}[htbp]
\centering
\small
\caption{Selected model-based direct variance-payoff strategies}
\label{tab:direct_main}
\begin{tabularx}{\textwidth}{Xrrrrrrrr}
\toprule
Market / strategy
& Return
& Vol.
& Sharpe
& Sortino
& MDD
& CVaR
& Turn.
& Obs. \\
\midrule
US High VRP
& 6.23\%
& 7.17\%
& 0.882
& 1.218
& -23.46\%
& -5.26\%
& 2.01\%
& 232 \\
EU VRP $>0$
& 13.08\%
& 9.70\%
& 1.324
& 1.891
& -21.36\%
& -7.34\%
& 4.39\%
& 170 \\
\bottomrule
\end{tabularx}
\end{table}

The United States High-VRP strategy generates an annualized return of 6.23\%, volatility of 7.17\%, and a Sharpe ratio of 0.882. Its maximum drawdown is -23.46\% and its empirical 95\% CVaR is -5.26\%.

On the same 232-observation direct-variance comparison sample, the 60/40 benchmark has a Sharpe ratio of 0.827 and equal weighting 0.852. The direct strategy therefore improves the Sharpe ratio on this particular aligned sample. This difference is not sufficient to establish investor-welfare dominance.

The alternative United States strategy active whenever lagged VRP is positive illustrates the importance of the exposure rule. It produces an annualized return of 8.18\%, but volatility rises to 16.07\%, the Sharpe ratio falls to 0.590, maximum drawdown reaches -39.77\%, and CVaR reaches -13.29\%. Selecting months with unusually high lagged VRP therefore produces a materially more defensive United States payoff profile than merely requiring VRP to be positive.

\subsection{European direct-variance evidence}

The European result is considerably stronger within the tested approximation. The selected strategy takes exposure when lagged VRP is positive. It produces an annualized return of 13.08\%, annualized volatility of 9.70\%, a Sharpe ratio of 1.324, and a Sortino ratio of 1.891. Maximum drawdown is -21.36\% and empirical 95\% CVaR is -7.34\%.

This result is economically different from the European allocation evidence. Adding VRP to HMM, RSM, or machine-learning allocation does not generate robust benchmark dominance. Yet mapping lagged implied variance into the model-based short-variance payoff produces substantially stronger risk-adjusted performance.

This is the first central cross-channel result of the thesis: weak evidence for VRP as a universal allocation feature does not imply weak variance carry. Information value and payoff value are distinct empirical objects.

The result must nevertheless retain the methodological qualification established earlier. The strategy is a model-based direct variance-payoff approximation. It is not an observed variance-swap return and does not include an exact option-replicated strike, contractual variance notional, collateral and margin processes, dealer bid--ask spreads, or daily mark-to-market.

Figure~\ref{fig:direct-variance-cross-market} summarizes the central cross-market direct-payoff result. Panel (a) compares the Sharpe ratio of the selected direct-variance strategy with the two balanced benchmarks. Panel (b) reports the corresponding fee-equivalent welfare difference at $\gamma=5$. The two panels highlight why Sharpe-ratio improvement alone is insufficient for the United States conclusion, while the European result remains economically much stronger under the welfare criterion.

\begin{figure}[htbp]
\centering

\begin{subfigure}[t]{0.49\textwidth}
\centering
\includegraphics[
    width=\linewidth
]{direct_variance_sharpe_cross_market.png}
\caption{Sharpe-ratio comparison}
\end{subfigure}
\hfill
\begin{subfigure}[t]{0.49\textwidth}
\centering
\includegraphics[
    width=\linewidth
]{direct_variance_welfare_cross_market.png}
\caption{Fee-equivalent welfare comparison}
\end{subfigure}

\caption{Cross-market economic comparison of the selected model-based direct variance-payoff strategies}
\label{fig:direct-variance-cross-market}
\end{figure}

\subsection{Investor-welfare evidence}

Risk-adjusted performance alone does not establish economic superiority, particularly for negatively skewed strategies. Table~\ref{tab:welfare_main} therefore reports the principal welfare comparison at risk aversion $\gamma=5$.

\begin{table}[htbp]
\centering
\small
\caption{Direct-variance welfare evidence at $\gamma=5$}
\label{tab:welfare_main}
\begin{tabularx}{\textwidth}{Xrrrrrr}
\toprule
Strategy
& MV CEQ
& CRRA CE
& vs 60/40
& CI low
& vs 1/N
& CI low \\
\midrule
US High VRP
& 5.04\%
& 5.00\%
& -70.9 bp
& -5.49\%
& -45.0 bp
& -4.94\% \\
EU VRP $>0$
& 10.49\%
& 10.67\%
& +926.6 bp
& +3.27\%
& +898.8 bp
& +2.91\% \\
\bottomrule
\end{tabularx}
\end{table}

The United States conclusion weakens once investor welfare is considered. The High-VRP strategy has an annualized mean--variance certainty equivalent of 5.04\% and a CRRA certainty equivalent of 5.00\%. Its mean--variance welfare difference is -70.9 basis points relative to 60/40 and -45.0 basis points relative to equal weighting. The bootstrap intervals include zero: the interval relative to 60/40 runs from approximately -5.49\% to +4.92\%, while the interval relative to equal weighting runs from -4.94\% to +4.59\%. The United States direct strategy therefore does not establish statistically robust welfare dominance.

The European result is materially different. The VRP-positive strategy has a mean--variance certainty equivalent of 10.49\% and a CRRA certainty equivalent of 10.67\%. Its mean--variance advantage is 926.6 basis points relative to 60/40 and 898.8 basis points relative to equal weighting. More importantly, the lower bootstrap bounds remain positive at +3.27\% and +2.91\%, respectively.

The European evidence therefore survives a more demanding economic criterion than the Sharpe ratio alone. Within the tested direct-payoff approximation, the selected European strategy produces positive investor-welfare differences against both traditional portfolios and those differences remain positive at the lower bootstrap bounds.

\subsection{Empirical synthesis}

Three findings emerge from the combined evidence.

First, the allocation channel is economically modest. In the United States, the best HMM, RSM, and machine-learning specifications are competitive in selected dimensions but do not dominate the traditional portfolios. The most visible benefit is sometimes lower drawdown, obtained at the cost of substantially higher turnover. In Europe, the evidence is weaker: the strongest dynamic models remain below the strongest benchmark Sharpe ratios.

Second, predictive information and investment value are not equivalent. European Random Forest results indicate that VRP variables contain incremental stress-classification information, yet this improvement does not generate statistically significant allocation welfare. The same distinction appears across the regime models, where adding VRP changes state probabilities and model rankings without producing robust portfolio dominance.

Third, the payoff channel produces the strongest cross-market asymmetry. The United States direct-variance strategy has attractive defensive characteristics but does not establish robust welfare superiority. In Europe, the selected direct variance-payoff approximation generates a Sharpe ratio of 1.324 and large positive certainty-equivalent differences with positive lower bootstrap bounds.

The results therefore do not support the statement that VRP is universally more useful as a state variable, nor do they support a claim that VRP fails in Europe. Instead, they imply that the economic value associated with variance information depends critically on the payoff structure through which it is harvested.

The next section tests whether these conclusions survive changes in implementation assumptions, risk-estimation windows, exposure limits, subperiods, crisis exclusions, rebalancing rules, and investor risk aversion.
